% --- Archon physics formalization source begin ---
% archon:physics
% archon:covers PhyXMiniProblems/problem_phyx_mini_0744.lean
% archon:source-report reports/phyx_mini/problem_phyx_mini_0744.source.json

\chapter{Physics problem phyx\_mini\_0744}
\label{ch:PhyXMiniProblems_problem_phyx_mini_0744}

\paragraph{Problem source.}
\#\# Physical scenario

In figure, a ball is launched with a velocity of magnitude \$10.0\textbackslash{} m/s\$, at an angle of \$50.0\textasciicircum{}\textbackslash{}circ\$ to the horizontal. The launch point is at the base of a ramp of horizontal length \$d\_1 = 6.00\textbackslash{} m\$ and height \$d\_2 = 3.60\textbackslash{} m\$. A plateau is located at the top of the ramp.

\#\# Question

When it lands, what is the angle of its displacement from the launch point?

\#\# Answer choices

- A: A: \$26.4\textasciicircum{}\textbackslash{}circ\$

- B: B: \$28.8\textasciicircum{}\textbackslash{}circ\$

- C: C: \$31.0\textasciicircum{}\textbackslash{}circ\$

- D: D: \$36.2\textasciicircum{}\textbackslash{}circ\$

\#\# Figure caption (auxiliary; use the image as primary evidence)

The image is a physics diagram illustrating the motion of a ball. It comprises:

- A small green ball labeled as "Ball."
- A magenta arrow labeled \textbackslash{}( \textbackslash{}vec\{v\_0\} \textbackslash{}) originating from the ball, indicating its initial velocity and direction.
- A ramp with two distinct sections: the first section is inclined upwards, and the second section is horizontal.
- Two distance labels:
  - \textbackslash{}( d\_1 \textbackslash{}) is the horizontal distance from the starting position of the ball to where the horizontal section begins.
  - \textbackslash{}( d\_2 \textbackslash{}) is the vertical distance representing the height of the horizontal section above the initial level.
- The background is shaded in light pink to highlight the path of the ramp.

\#\# Dataset metadata

Kinematics; Physical Model Grounding Reasoning, Spatial Relation Reasoning

\paragraph{Recorded answer/context.}
C: C: \$31.0\textasciicircum{}\textbackslash{}circ\$

\paragraph{Figure/image path.}
/root/proposal\_for\_physic/science-mango/phyx\_mini\_run/phyx\_data/test\_image/744.png

\paragraph{Formalization target.}
create a compiling Lean file with sorry bodies at `PhyXMiniProblems/problem\_phyx\_mini\_0744.lean`. The Lean declarations must preserve the physical quantities, dimensions or dimensional roles, figure labels, governing-law hypotheses, and final relation expressed by this problem.
Use Mathlib/Physlib names found through LeanExplore where available. If a physics API is missing, introduce faithful local abstractions rather than scalar placeholder aliases.

\begin{theorem}[Physics formalization target]
\label{thm:physics:phyx_mini_0744:target}
\lean{PhyXMiniProblems.ProblemPhyXMini0744.landingDisplacementAngleIsAnswerC}
\uses{def:physics:phyx-mini-0744:phyxminiproblems-problemphyxmini0744-ballrampprojectilesetup, def:physics:phyx-mini-0744:phyxminiproblems-problemphyxmini0744-matchesproblemandprimaryfigure, def:physics:phyx-mini-0744:phyxminiproblems-problemphyxmini0744-isoninclinedramp, def:physics:phyx-mini-0744:phyxminiproblems-problemphyxmini0744-haspositivephysicalparameters, def:physics:phyx-mini-0744:phyxminiproblems-problemphyxmini0744-usesstandardearthgravity, def:physics:phyx-mini-0744:phyxminiproblems-problemphyxmini0744-obeysuniformgravityprojectilelaw, def:physics:phyx-mini-0744:phyxminiproblems-problemphyxmini0744-obeysfirstterraincontactlaw, def:physics:phyx-mini-0744:phyxminiproblems-problemphyxmini0744-landingdisplacementangleradians, def:physics:phyx-mini-0744:phyxminiproblems-problemphyxmini0744-landingdisplacementangledegrees, def:physics:phyx-mini-0744:phyxminiproblems-problemphyxmini0744-answerchoice, def:physics:phyx-mini-0744:phyxminiproblems-problemphyxmini0744-roundstonearesttenth, def:physics:phyx-mini-0744:phyxminiproblems-problemphyxmini0744-isnearestanswerchoice}
The assigned autoformalize agent should translate this physics problem into Lean declarations in the covered file, with theorem and lemma proof bodies written as `by sorry`.
\end{theorem}
\begin{proof}
This is an autoformalization task, not a proof task. Produce statements that can later be proved by the physics prover without weakening the physical model.
\end{proof}

% --- Archon declaration topology begin ---
\section{Declaration topology}
\begin{definition}[Length Quantity]
  \label{def:physics:phyx-mini-0744:phyxminiproblems-problemphyxmini0744-lengthquantity}
  \lean{PhyXMiniProblems.ProblemPhyXMini0744.LengthQuantity}
  A nonnegative physical length
\end{definition}
\begin{proof}
  This is the declared model definition of Length Quantity.
\end{proof}
\begin{definition}[Time Quantity]
  \label{def:physics:phyx-mini-0744:phyxminiproblems-problemphyxmini0744-timequantity}
  \lean{PhyXMiniProblems.ProblemPhyXMini0744.TimeQuantity}
  A nonnegative elapsed time
\end{definition}
\begin{proof}
  This is the declared model definition of Time Quantity.
\end{proof}
\begin{definition}[Speed Quantity]
  \label{def:physics:phyx-mini-0744:phyxminiproblems-problemphyxmini0744-speedquantity}
  \lean{PhyXMiniProblems.ProblemPhyXMini0744.SpeedQuantity}
  A nonnegative physical speed
\end{definition}
\begin{proof}
  This is the declared model definition of Speed Quantity.
\end{proof}
\begin{definition}[Acceleration Quantity]
  \label{def:physics:phyx-mini-0744:phyxminiproblems-problemphyxmini0744-accelerationquantity}
  \lean{PhyXMiniProblems.ProblemPhyXMini0744.AccelerationQuantity}
  A nonnegative acceleration magnitude
\end{definition}
\begin{proof}
  This is the declared model definition of Acceleration Quantity.
\end{proof}
\begin{definition}[quantity Readout]
  \label{def:physics:phyx-mini-0744:phyxminiproblems-problemphyxmini0744-quantityreadout}
  \lean{PhyXMiniProblems.ProblemPhyXMini0744.quantityReadout}
  Read a nonnegative dimensionful quantity in coherent SI units
\end{definition}
\begin{proof}
  This is the declared model definition of quantity Readout.
\end{proof}
\begin{definition}[length In Meters]
  \label{def:physics:phyx-mini-0744:phyxminiproblems-problemphyxmini0744-lengthinmeters}
  \lean{PhyXMiniProblems.ProblemPhyXMini0744.lengthInMeters}
  \uses{def:physics:phyx-mini-0744:phyxminiproblems-problemphyxmini0744-lengthquantity, def:physics:phyx-mini-0744:phyxminiproblems-problemphyxmini0744-quantityreadout}
  Metre readout of a physical length
\end{definition}
\begin{proof}
  This is the declared model definition of length In Meters.
\end{proof}
\begin{definition}[time In Seconds]
  \label{def:physics:phyx-mini-0744:phyxminiproblems-problemphyxmini0744-timeinseconds}
  \lean{PhyXMiniProblems.ProblemPhyXMini0744.timeInSeconds}
  \uses{def:physics:phyx-mini-0744:phyxminiproblems-problemphyxmini0744-timequantity, def:physics:phyx-mini-0744:phyxminiproblems-problemphyxmini0744-quantityreadout}
  Second readout of a physical elapsed time
\end{definition}
\begin{proof}
  This is the declared model definition of time In Seconds.
\end{proof}
\begin{definition}[speed In Meters Per Second]
  \label{def:physics:phyx-mini-0744:phyxminiproblems-problemphyxmini0744-speedinmeterspersecond}
  \lean{PhyXMiniProblems.ProblemPhyXMini0744.speedInMetersPerSecond}
  \uses{def:physics:phyx-mini-0744:phyxminiproblems-problemphyxmini0744-speedquantity, def:physics:phyx-mini-0744:phyxminiproblems-problemphyxmini0744-quantityreadout}
  Metres-per-second readout of a physical speed
\end{definition}
\begin{proof}
  This is the declared model definition of speed In Meters Per Second.
\end{proof}
\begin{definition}[acceleration In Meters Per Second Squared]
  \label{def:physics:phyx-mini-0744:phyxminiproblems-problemphyxmini0744-accelerationinmeterspersecondsquared}
  \lean{PhyXMiniProblems.ProblemPhyXMini0744.accelerationInMetersPerSecondSquared}
  \uses{def:physics:phyx-mini-0744:phyxminiproblems-problemphyxmini0744-accelerationquantity, def:physics:phyx-mini-0744:phyxminiproblems-problemphyxmini0744-quantityreadout}
  Metres-per-second-squared readout of an acceleration magnitude
\end{definition}
\begin{proof}
  This is the declared model definition of acceleration In Meters Per Second Squared.
\end{proof}
\begin{definition}[angle From Degrees]
  \label{def:physics:phyx-mini-0744:phyxminiproblems-problemphyxmini0744-anglefromdegrees}
  \lean{PhyXMiniProblems.ProblemPhyXMini0744.angleFromDegrees}
  Convert a degree readout to Mathlib's physical angle type
\end{definition}
\begin{proof}
  This is the declared model definition of angle From Degrees.
\end{proof}
\begin{definition}[Physical Point2 D]
  \label{def:physics:phyx-mini-0744:phyxminiproblems-problemphyxmini0744-physicalpoint2d}
  \lean{PhyXMiniProblems.ProblemPhyXMini0744.PhysicalPoint2D}
  \uses{def:physics:phyx-mini-0744:phyxminiproblems-problemphyxmini0744-lengthquantity}
  A point in the vertical cross-section, with two physical length coordinates measured from a common Cartesian origin
\end{definition}
\begin{proof}
  This is the declared model definition of Physical Point2 D.
\end{proof}
\begin{definition}[x Meters]
  \label{def:physics:phyx-mini-0744:phyxminiproblems-problemphyxmini0744-physicalpoint2d-xmeters}
  \lean{PhyXMiniProblems.ProblemPhyXMini0744.PhysicalPoint2D.xMeters}
  \uses{def:physics:phyx-mini-0744:phyxminiproblems-problemphyxmini0744-lengthinmeters, def:physics:phyx-mini-0744:phyxminiproblems-problemphyxmini0744-physicalpoint2d}
  Horizontal coordinate of a physical point, in metres
\end{definition}
\begin{proof}
  This is the declared model definition of x Meters.
\end{proof}
\begin{definition}[y Meters]
  \label{def:physics:phyx-mini-0744:phyxminiproblems-problemphyxmini0744-physicalpoint2d-ymeters}
  \lean{PhyXMiniProblems.ProblemPhyXMini0744.PhysicalPoint2D.yMeters}
  \uses{def:physics:phyx-mini-0744:phyxminiproblems-problemphyxmini0744-lengthinmeters, def:physics:phyx-mini-0744:phyxminiproblems-problemphyxmini0744-physicalpoint2d}
  Vertical coordinate of a physical point, in metres
\end{definition}
\begin{proof}
  This is the declared model definition of y Meters.
\end{proof}
\begin{definition}[Figure Component]
  \label{def:physics:phyx-mini-0744:phyxminiproblems-problemphyxmini0744-figurecomponent}
  \lean{PhyXMiniProblems.ProblemPhyXMini0744.FigureComponent}
  Individually identifiable objects and labels visible in the primary figure
\end{definition}
\begin{proof}
  This is the declared model definition of Figure Component.
\end{proof}
\begin{definition}[Terrain Section]
  \label{def:physics:phyx-mini-0744:phyxminiproblems-problemphyxmini0744-terrainsection}
  \lean{PhyXMiniProblems.ProblemPhyXMini0744.TerrainSection}
  The two pieces of the terrain on which the ball could first land
\end{definition}
\begin{proof}
  This is the declared model definition of Terrain Section.
\end{proof}
\begin{definition}[Ball Ramp Projectile Setup]
  \label{def:physics:phyx-mini-0744:phyxminiproblems-problemphyxmini0744-ballrampprojectilesetup}
  \lean{PhyXMiniProblems.ProblemPhyXMini0744.BallRampProjectileSetup}
  \uses{def:physics:phyx-mini-0744:phyxminiproblems-problemphyxmini0744-lengthquantity, def:physics:phyx-mini-0744:phyxminiproblems-problemphyxmini0744-timequantity, def:physics:phyx-mini-0744:phyxminiproblems-problemphyxmini0744-speedquantity, def:physics:phyx-mini-0744:phyxminiproblems-problemphyxmini0744-accelerationquantity, def:physics:phyx-mini-0744:phyxminiproblems-problemphyxmini0744-physicalpoint2d, def:physics:phyx-mini-0744:phyxminiproblems-problemphyxmini0744-figurecomponent}
  The independent physical objects and quantities in the problem. Ball is kept abstract so that the ball is not collapsed to scalar data. In particular, the landing location, landing time, and trajectory are not fixed to the requested answer
\end{definition}
\begin{proof}
  This is the declared model definition of Ball Ramp Projectile Setup.
\end{proof}
\begin{definition}[Matches Problem And Primary Figure]
  \label{def:physics:phyx-mini-0744:phyxminiproblems-problemphyxmini0744-matchesproblemandprimaryfigure}
  \lean{PhyXMiniProblems.ProblemPhyXMini0744.MatchesProblemAndPrimaryFigure}
  \uses{def:physics:phyx-mini-0744:phyxminiproblems-problemphyxmini0744-lengthinmeters, def:physics:phyx-mini-0744:phyxminiproblems-problemphyxmini0744-speedinmeterspersecond, def:physics:phyx-mini-0744:phyxminiproblems-problemphyxmini0744-anglefromdegrees, def:physics:phyx-mini-0744:phyxminiproblems-problemphyxmini0744-ballrampprojectilesetup}
  Exact problem-statement and primary-image readouts. This records the straight ramp's endpoint geometry and the launch data, but neither selects a terrain section for the landing nor assigns the displacement angle
\end{definition}
\begin{proof}
  This is the declared model definition of Matches Problem And Primary Figure.
\end{proof}
\begin{definition}[Is On Inclined Ramp]
  \label{def:physics:phyx-mini-0744:phyxminiproblems-problemphyxmini0744-isoninclinedramp}
  \lean{PhyXMiniProblems.ProblemPhyXMini0744.IsOnInclinedRamp}
  \uses{def:physics:phyx-mini-0744:phyxminiproblems-problemphyxmini0744-lengthinmeters, def:physics:phyx-mini-0744:phyxminiproblems-problemphyxmini0744-physicalpoint2d, def:physics:phyx-mini-0744:phyxminiproblems-problemphyxmini0744-ballrampprojectilesetup}
  A point lies on the closed straight incline joining the launch point to the ramp top. The last equation is dimensionally a metre-squared equality; its quotient form would say dy / dx = d₂ / d₁
\end{definition}
\begin{proof}
  This is the declared model definition of Is On Inclined Ramp.
\end{proof}
\begin{definition}[Is On Plateau]
  \label{def:physics:phyx-mini-0744:phyxminiproblems-problemphyxmini0744-isonplateau}
  \lean{PhyXMiniProblems.ProblemPhyXMini0744.IsOnPlateau}
  \uses{def:physics:phyx-mini-0744:phyxminiproblems-problemphyxmini0744-lengthinmeters, def:physics:phyx-mini-0744:phyxminiproblems-problemphyxmini0744-physicalpoint2d, def:physics:phyx-mini-0744:phyxminiproblems-problemphyxmini0744-ballrampprojectilesetup}
  A point lies on the horizontal plateau beginning at the top of the ramp
\end{definition}
\begin{proof}
  This is the declared model definition of Is On Plateau.
\end{proof}
\begin{definition}[Is On Terrain Surface]
  \label{def:physics:phyx-mini-0744:phyxminiproblems-problemphyxmini0744-isonterrainsurface}
  \lean{PhyXMiniProblems.ProblemPhyXMini0744.IsOnTerrainSurface}
  \uses{def:physics:phyx-mini-0744:phyxminiproblems-problemphyxmini0744-physicalpoint2d, def:physics:phyx-mini-0744:phyxminiproblems-problemphyxmini0744-ballrampprojectilesetup, def:physics:phyx-mini-0744:phyxminiproblems-problemphyxmini0744-isoninclinedramp, def:physics:phyx-mini-0744:phyxminiproblems-problemphyxmini0744-isonplateau}
  A point lies on one of the two exposed terrain sections in the diagram
\end{definition}
\begin{proof}
  This is the declared model definition of Is On Terrain Surface.
\end{proof}
\begin{definition}[Has Positive Physical Parameters]
  \label{def:physics:phyx-mini-0744:phyxminiproblems-problemphyxmini0744-haspositivephysicalparameters}
  \lean{PhyXMiniProblems.ProblemPhyXMini0744.HasPositivePhysicalParameters}
  \uses{def:physics:phyx-mini-0744:phyxminiproblems-problemphyxmini0744-lengthinmeters, def:physics:phyx-mini-0744:phyxminiproblems-problemphyxmini0744-timeinseconds, def:physics:phyx-mini-0744:phyxminiproblems-problemphyxmini0744-speedinmeterspersecond, def:physics:phyx-mini-0744:phyxminiproblems-problemphyxmini0744-accelerationinmeterspersecondsquared, def:physics:phyx-mini-0744:phyxminiproblems-problemphyxmini0744-ballrampprojectilesetup}
  Positivity and forward-motion conditions for the physical branch depicted in the image
\end{definition}
\begin{proof}
  This is the declared model definition of Has Positive Physical Parameters.
\end{proof}
\begin{definition}[Uses Standard Earth Gravity]
  \label{def:physics:phyx-mini-0744:phyxminiproblems-problemphyxmini0744-usesstandardearthgravity}
  \lean{PhyXMiniProblems.ProblemPhyXMini0744.UsesStandardEarthGravity}
  \uses{def:physics:phyx-mini-0744:phyxminiproblems-problemphyxmini0744-accelerationinmeterspersecondsquared, def:physics:phyx-mini-0744:phyxminiproblems-problemphyxmini0744-ballrampprojectilesetup}
  The conventional near-Earth gravitational acceleration used by the elementary projectile model
\end{definition}
\begin{proof}
  This is the declared model definition of Uses Standard Earth Gravity.
\end{proof}
\begin{definition}[Obeys Uniform Gravity Projectile Law]
  \label{def:physics:phyx-mini-0744:phyxminiproblems-problemphyxmini0744-obeysuniformgravityprojectilelaw}
  \lean{PhyXMiniProblems.ProblemPhyXMini0744.ObeysUniformGravityProjectileLaw}
  \uses{def:physics:phyx-mini-0744:phyxminiproblems-problemphyxmini0744-timeinseconds, def:physics:phyx-mini-0744:phyxminiproblems-problemphyxmini0744-speedinmeterspersecond, def:physics:phyx-mini-0744:phyxminiproblems-problemphyxmini0744-accelerationinmeterspersecondsquared, def:physics:phyx-mini-0744:phyxminiproblems-problemphyxmini0744-ballrampprojectilesetup}
  The standard no-drag projectile equations under constant downward gravity. They are imposed only through the stated landing time, so the nonnegative physical coordinate representation is not required to describe the later mathematical continuation below the terrain
\end{definition}
\begin{proof}
  This is the declared model definition of Obeys Uniform Gravity Projectile Law.
\end{proof}
\begin{definition}[Obeys First Terrain Contact Law]
  \label{def:physics:phyx-mini-0744:phyxminiproblems-problemphyxmini0744-obeysfirstterraincontactlaw}
  \lean{PhyXMiniProblems.ProblemPhyXMini0744.ObeysFirstTerrainContactLaw}
  \uses{def:physics:phyx-mini-0744:phyxminiproblems-problemphyxmini0744-timeinseconds, def:physics:phyx-mini-0744:phyxminiproblems-problemphyxmini0744-ballrampprojectilesetup, def:physics:phyx-mini-0744:phyxminiproblems-problemphyxmini0744-isonterrainsurface}
  Landing means the first positive-time contact of the ballistic trajectory with either the incline or the plateau. The disjunction deliberately does not assume which terrain section is reached; distinguishing the two is part of the physics argument
\end{definition}
\begin{proof}
  This is the declared model definition of Obeys First Terrain Contact Law.
\end{proof}
\begin{definition}[landing Displacement Angle Radians]
  \label{def:physics:phyx-mini-0744:phyxminiproblems-problemphyxmini0744-landingdisplacementangleradians}
  \lean{PhyXMiniProblems.ProblemPhyXMini0744.landingDisplacementAngleRadians}
  \uses{def:physics:phyx-mini-0744:phyxminiproblems-problemphyxmini0744-ballrampprojectilesetup}
  The principal radian angle of the launch-to-landing displacement from the positive horizontal direction. Forward motion makes Real.arctan (dy / dx) the appropriate branch
\end{definition}
\begin{proof}
  This is the declared model definition of landing Displacement Angle Radians.
\end{proof}
\begin{definition}[landing Displacement Angle Degrees]
  \label{def:physics:phyx-mini-0744:phyxminiproblems-problemphyxmini0744-landingdisplacementangledegrees}
  \lean{PhyXMiniProblems.ProblemPhyXMini0744.landingDisplacementAngleDegrees}
  \uses{def:physics:phyx-mini-0744:phyxminiproblems-problemphyxmini0744-ballrampprojectilesetup, def:physics:phyx-mini-0744:phyxminiproblems-problemphyxmini0744-landingdisplacementangleradians}
  Degree readout of the launch-to-landing displacement angle
\end{definition}
\begin{proof}
  This is the declared model definition of landing Displacement Angle Degrees.
\end{proof}
\begin{definition}[Answer Choice]
  \label{def:physics:phyx-mini-0744:phyxminiproblems-problemphyxmini0744-answerchoice}
  \lean{PhyXMiniProblems.ProblemPhyXMini0744.AnswerChoice}
  Labels printed beside the four displayed answer choices
\end{definition}
\begin{proof}
  This is the declared model definition of Answer Choice.
\end{proof}
\begin{definition}[angle Degrees]
  \label{def:physics:phyx-mini-0744:phyxminiproblems-problemphyxmini0744-answerchoice-angledegrees}
  \lean{PhyXMiniProblems.ProblemPhyXMini0744.AnswerChoice.angleDegrees}
  \uses{def:physics:phyx-mini-0744:phyxminiproblems-problemphyxmini0744-answerchoice}
  Degree readout printed for each answer choice
\end{definition}
\begin{proof}
  This is the declared model definition of angle Degrees.
\end{proof}
\begin{definition}[recorded Answer Choice]
  \label{def:physics:phyx-mini-0744:phyxminiproblems-problemphyxmini0744-recordedanswerchoice}
  \lean{PhyXMiniProblems.ProblemPhyXMini0744.recordedAnswerChoice}
  \uses{def:physics:phyx-mini-0744:phyxminiproblems-problemphyxmini0744-answerchoice}
  Dataset metadata: the recorded answer label is C. This definition is not used as a premise of either theorem
\end{definition}
\begin{proof}
  This is the declared model definition of recorded Answer Choice.
\end{proof}
\begin{definition}[Rounds To Nearest Tenth]
  \label{def:physics:phyx-mini-0744:phyxminiproblems-problemphyxmini0744-roundstonearesttenth}
  \lean{PhyXMiniProblems.ProblemPhyXMini0744.RoundsToNearestTenth}
  The exact angle rounds to a displayed value to the nearest tenth of a degree
\end{definition}
\begin{proof}
  This is the declared model definition of Rounds To Nearest Tenth.
\end{proof}
\begin{definition}[Is Nearest Answer Choice]
  \label{def:physics:phyx-mini-0744:phyxminiproblems-problemphyxmini0744-isnearestanswerchoice}
  \lean{PhyXMiniProblems.ProblemPhyXMini0744.IsNearestAnswerChoice}
  \uses{def:physics:phyx-mini-0744:phyxminiproblems-problemphyxmini0744-answerchoice}
  A displayed answer is strictly closer to the physical angle than every other displayed answer
\end{definition}
\begin{proof}
  This is the declared model definition of Is Nearest Answer Choice.
\end{proof}
\begin{lemma}[landing Point Lies On Inclined Ramp]
  \label{lem:physics:phyx-mini-0744:phyxminiproblems-problemphyxmini0744-landingpointliesoninclinedramp}
  \lean{PhyXMiniProblems.ProblemPhyXMini0744.landingPointLiesOnInclinedRamp}
  \uses{def:physics:phyx-mini-0744:phyxminiproblems-problemphyxmini0744-ballrampprojectilesetup, def:physics:phyx-mini-0744:phyxminiproblems-problemphyxmini0744-matchesproblemandprimaryfigure, def:physics:phyx-mini-0744:phyxminiproblems-problemphyxmini0744-isoninclinedramp, def:physics:phyx-mini-0744:phyxminiproblems-problemphyxmini0744-haspositivephysicalparameters, def:physics:phyx-mini-0744:phyxminiproblems-problemphyxmini0744-usesstandardearthgravity, def:physics:phyx-mini-0744:phyxminiproblems-problemphyxmini0744-obeysuniformgravityprojectilelaw, def:physics:phyx-mini-0744:phyxminiproblems-problemphyxmini0744-obeysfirstterraincontactlaw}
  The 10 m/s, 50° trajectory under standard gravity cannot reach the 3.60 m plateau; its first positive terrain contact is therefore on the inclined section before d₁ = 6.00 m
\end{lemma}
\begin{proof}
  The conclusion follows from the governing relations and figure readouts named in the statement of landing Point Lies On Inclined Ramp.
\end{proof}
% --- Archon declaration topology end ---
% --- Archon physics formalization source end ---
